\section{Introduction}
\label{sec:introduction}

Modern data systems increasingly embed hybrid and vector search. Production hybrid-search stacks combine BM25~\cite{robertson2009bm25}, dense vector search~\cite{wang2022e5}, learned-sparse retrieval~\cite{formal2021splade}, late interaction~\cite{santhanam2021colbertv2}, reciprocal-rank fusion~\cite{cormack2009rrf}, metadata-filter placement, and neural reranking~\cite{pradeep2023rankzephyr}. This is the operator menu that the retrieval literature and current vector engines expose in general; we do not model one particular deployment, and the action catalog we actually execute is listed in Table~\ref{tab:catalog}. These are \emph{physical plans} with different latency profiles, filter behavior, and failure modes, so the systems question is a query-optimization one: which physical plan should serve a workload under a latency SLO and a filter constraint, and, critically, can that choice be \emph{trusted} once the workload drifts from what was seen at calibration. A plan that is fast on yesterday's queries can silently start violating the latency SLO when today's workload shifts toward harder, higher-selectivity queries. In deployment the quantity that matters is the \emph{rate} at which the chosen plan breaches the SLO over a serving window: this rate drives SLA penalties and admission decisions, yet current selectors set it by hoping a plan generalizes, with no guarantee before execution. Bounding it ahead of time---so an operator can audit, and admit or hold back, a plan before it serves traffic---is the gap we close.

Most current strategies choose at the wrong unit and offer no such trust. Fixed recipes such as BM25-only, dense-only, fixed fusion, or fixed reranking depth are simple but brittle: the best plan changes with selectivity, budget, and query mix. Per-query routers are more adaptive, but they decide from isolated query features, and the documented unreliability of query performance prediction makes per-query selection fragile in exactly the constrained regime we care about. Neither family, however, says anything about \emph{how often the deployed choice will violate the SLO on the next, possibly shifted, workload window}. An arriving workload window, however, can be observed before execution: it reveals stable regions where the same plan is repeatedly good, and it lets a system bound its own risk instead of merely hoping the chosen plan generalizes.

This paper treats hybrid-search plan selection as \emph{workload-window physical-plan compilation with an auditable deployment guarantee}. The input is a labeled training workload, an unlabeled evaluation workload window, and telemetry from candidate operators; the output is a compiled design---queries partitioned into plan cells, a physical plan bound to each cell, a cell-level plan frontier---together with a \emph{distribution-free, finite-sample certificate} on the eval-window SLO-violation rate that holds under a bounded amount of workload shift. The certificate is operational: when a cell's plan is not certifiably safe, the compiler abstains to a certifiably feasible fallback. The operators themselves are unchanged; what is new is that the optimizer certifies \emph{which plan to deploy, and when to abstain}, the way a database physical-design tool reasons about a workload rather than a single query. To our knowledge no prior hybrid-search method provides such a guarantee, and none of our baselines provides one as used.

Our system is \CWC{}, the \emph{Certified Workload-window Compiler}. \CWC{} first embeds train and evaluation queries using observable workload telemetry---cheap retrieval probes, score-shape features, branch overlap, filter survival, and latency signals---then partitions the workload into plan cells and binds each cell to the action with the best training constrained utility, so each evaluation query inherits the plan of its cell. On top of this it adds an abstain rule and a deployment certificate, both developed later in the paper. We write \WCR{} for the \emph{un-guarded} variant (the same compiler without the abstain rule); it is used only as an ablation and in the preliminary single-collection pilot and cross-dataset audit, and when no cell abstains \CWC{}$\equiv$\WCR{}. The method is transductive in the systems sense: it uses unlabeled telemetry from the current workload window, but never evaluation relevance labels.

Our canonical evidence (Section~\ref{sec:results-rank}) is a $45$-block study over SciFact, FiQA, and NFCorpus with real external operators under a controlled bounded workload shift; a single-collection NFCorpus pilot and a $360$-block cross-dataset audit corroborate it as robustness, not the headline.

The claim is intentionally scoped: optimizer-level gains plus an auditable deployment guarantee under a fixed action catalog, not universal superiority over every retriever stack. In particular \CWC{} \emph{ties} the strongest calibration-aware selectors on mean utility; its differentiator there is the certificate, not raw rank.

This paper makes three contributions.
\begin{enumerate}
  \item \textbf{Abstraction.} We formulate hybrid-search plan selection as \emph{certified workload-window physical-plan compilation}: an unlabeled workload window is partitioned into plan cells, each cell is bound to an executable physical plan, and each deployment decision carries an auditable risk guarantee.
  \item \textbf{Guarantee.} We prove a distribution-free, finite-sample certificate (Proposition~\ref{prop:cert}) that upper-bounds the eval-window SLO-violation rate under a bounded within-cell shift budget, made operational through \CWC{}, a guarded compiler that abstains to a feasibility-best fallback on fragile cells; a density-ratio upper-confidence-bound variant removes the assumed-shift caveat. To our knowledge this is the first such guarantee at the \emph{workload-window} granularity for hybrid-search plan selection.
  \item \textbf{Evaluation.} On a $45$-block SciFact+FiQA+NFCorpus workload with BM25, E5, Qwen3-Embedding, SPLADE, ColBERTv2, fusion, and reranking actions, \CWC{} is average-rank first on constrained utility, Holm-beats four strong static operators, and ties calibration-aware selectors while adding a certificate they lack, cutting realized SLO violations from $\approx11\%$ to $\approx0.1\%$ at $\approx4\%$ nDCG cost; under an SLA admission-control objective it separates from those selectors (rank $1.44$, Holm-beating all six baselines), and a scaling study confirms the bound tightens with per-cell calibration size. Ablations and a $360$-block audit isolate the source of the gains.
\end{enumerate}
